\documentclass[aspectratio=169,11pt]{beamer}
\usepackage{../../../shared/theme/esmad_beamer_theme}
\usepackage{amsmath,amssymb}
\usepackage{booktabs}

\title[Regression Discontinuity]{Regression Discontinuity Designs}
\subtitle{Local Causal Identification at a Threshold}
\date{\today}

\newcommand{\E}{\mathbb{E}}

\begin{document}

\begin{frame}
\titlepage
\end{frame}

\begin{frame}{Learning Goals}
\begin{itemize}
\item explain why a treatment threshold can create quasi-experimental variation;
\item distinguish sharp and fuzzy regression discontinuity designs;
\item define the local causal estimand at the cutoff;
\item understand bandwidth choice and local polynomial estimation;
\item diagnose manipulation, sorting, and covariate discontinuities;
\item interpret robust bias-corrected uncertainty and limits on external validity.
\end{itemize}
\end{frame}

\begin{frame}{Outline}
\tableofcontents
\end{frame}

\section{The Basic Design}

\begin{frame}{Treatment Assignment at a Cutoff}
Let \(X_i\) be a running variable and \(c\) a known cutoff.

In a sharp RDD:
\[
D_i=\mathbb{1}\{X_i\ge c\}.
\]

Units just below and just above \(c\) are compared.

\begin{alertblock}{Identification intuition}
Near the threshold, units may be very similar except for the treatment status induced by the assignment rule.
\end{alertblock}
\end{frame}

\begin{frame}{The Local Causal Estimand}
The target is the treatment effect at the cutoff:
\[
\tau_{RD}
=
\E[Y_i(1)-Y_i(0)\mid X_i=c].
\]

This is a local estimand.

\begin{block}{Implication}
RDD does not automatically identify an average effect for the whole population.
\end{block}
\end{frame}

\section{Continuity Assumptions}

\begin{frame}{Continuity of Potential Outcomes}
A core assumption is that
\[
\E[Y_i(0)\mid X_i=x]
\]
and
\[
\E[Y_i(1)\mid X_i=x]
\]
are continuous at \(x=c\).

Then any jump in the observed conditional mean at the cutoff can be attributed to treatment.
\end{frame}

\begin{frame}{Why Continuity Can Be Plausible}
Continuity is plausible when:
\begin{itemize}
\item the cutoff is administratively imposed;
\item units cannot precisely manipulate the running variable;
\item other determinants of the outcome evolve smoothly near the threshold.
\end{itemize}
\end{frame}

\section{Sharp RDD}

\begin{frame}{Sharp RDD Estimand}
With deterministic assignment:
\[
\tau_{SRD}
=
\lim_{x\downarrow c}\E[Y\mid X=x]
-
\lim_{x\uparrow c}\E[Y\mid X=x].
\]

The discontinuity in outcomes identifies the local treatment effect.
\end{frame}

\begin{frame}{Local Linear Regression}
Estimate separate local lines on each side:
\[
Y_i
=
\alpha
+
\tau D_i
+
\beta_1(X_i-c)
+
\beta_2D_i(X_i-c)
+
u_i.
\]

Only observations within a bandwidth around \(c\) are used.
\end{frame}

\section{Bandwidth Choice}

\begin{frame}{Bias--Variance Tradeoff}
A narrow bandwidth:
\begin{itemize}
\item improves local comparability;
\item reduces functional-form dependence;
\item increases variance.
\end{itemize}

A wide bandwidth:
\begin{itemize}
\item uses more observations;
\item reduces variance;
\item risks bias from misspecified trends away from the cutoff.
\end{itemize}
\end{frame}

\begin{frame}{Bandwidth Sensitivity}
Report estimates across reasonable bandwidths.

\begin{alertblock}{Caution}
A result that appears only for one narrow bandwidth may be fragile.
\end{alertblock}
\end{frame}

\section{Polynomial Pitfalls}

\begin{frame}{Why High-Order Global Polynomials Are Risky}
High-order polynomial fits can:
\begin{itemize}
\item behave erratically near boundaries;
\item overfit distant observations;
\item create artificial curvature;
\item produce unstable extrapolation at the cutoff.
\end{itemize}

\begin{block}{Preferred approach}
Local linear or low-order local polynomial estimation is usually more stable.
\end{block}
\end{frame}

\section{Manipulation and Sorting}

\begin{frame}{Precise Manipulation Breaks the Design}
Suppose units can influence \(X_i\) to cross the threshold.

Then units just above and below \(c\) may differ systematically.

Examples:
\begin{itemize}
\item firms strategically reporting values;
\item students gaming test-score thresholds;
\item applicants manipulating eligibility metrics.
\end{itemize}
\end{frame}

\begin{frame}{Density Discontinuity}
A visible jump in the density of \(X\) at \(c\) can suggest sorting or manipulation.

\begin{alertblock}{Interpretation}
A density discontinuity is evidence against smooth assignment near the threshold, though context matters.
\end{alertblock}
\end{frame}

\section{Covariate Balance}

\begin{frame}{Predetermined Covariates Should Be Smooth}
Pre-treatment covariates should not jump at the cutoff.

Examples:
\begin{itemize}
\item age;
\item prior outcomes;
\item baseline income;
\item historical performance.
\end{itemize}

A discontinuity in predetermined covariates can signal invalid local comparability.
\end{frame}

\section{Placebo Cutoffs}

\begin{frame}{False Thresholds}
Estimate RDD effects at placebo cutoffs where treatment does not change.

Large placebo discontinuities may reveal:
\begin{itemize}
\item functional-form misspecification;
\item unrelated structural breaks;
\item periodicity or heaping in the running variable.
\end{itemize}
\end{frame}

\section{Fuzzy RDD}

\begin{frame}{When Treatment Probability Jumps, Not Treatment Itself}
In fuzzy RDD,
\[
\lim_{x\downarrow c}P(D=1\mid X=x)
-
\lim_{x\uparrow c}P(D=1\mid X=x)
\neq 0,
\]
but the jump is less than one.

The cutoff acts as an instrument for actual treatment.
\end{frame}

\begin{frame}{Fuzzy RDD as Local IV}
The fuzzy RDD estimand is
\[
\tau_{FRD}
=
\frac{
\lim_{x\downarrow c}\E[Y\mid X=x]
-
\lim_{x\uparrow c}\E[Y\mid X=x]
}{
\lim_{x\downarrow c}\E[D\mid X=x]
-
\lim_{x\uparrow c}\E[D\mid X=x]
}.
\]

This is a local Wald ratio.
\end{frame}

\begin{frame}{Interpretation of Fuzzy RDD}
Under monotonicity and standard assumptions, the fuzzy RDD identifies a local average treatment effect for units whose treatment status is changed by crossing the cutoff.
\end{frame}

\section{Inference}

\begin{frame}{Bias-Corrected Inference}
Local polynomial estimators have smoothing bias.

Modern RDD inference often uses:
\begin{itemize}
\item bias correction;
\item robust standard errors;
\item bandwidth selection tailored to inference.
\end{itemize}

\begin{block}{Reason}
Point-estimation bandwidths are not always optimal for confidence intervals.
\end{block}
\end{frame}

\section{Applied Reasoning}

\begin{frame}{Example: Eligibility Threshold}
Suppose a subsidy is available only to firms with score
\[
X_i\ge 70.
\]

To estimate the local effect near 70:
\begin{itemize}
\item plot outcomes against the running variable;
\item estimate local fits on each side;
\item test for score manipulation;
\item inspect baseline covariates;
\item vary bandwidths;
\item estimate placebo cutoffs.
\end{itemize}
\end{frame}

\begin{frame}{A Credible RDD Workflow}
\begin{enumerate}
\item Define the running variable and cutoff.
\item Explain the assignment rule.
\item State the local estimand.
\item Plot raw data near the cutoff.
\item Choose a principled bandwidth.
\item Estimate local linear models.
\item Test density continuity and covariate balance.
\item Run placebo cutoffs and bandwidth sensitivity.
\item Report robust bias-corrected uncertainty.
\item Limit interpretation to the population near the cutoff.
\end{enumerate}
\end{frame}

\begin{frame}{Common Failure Modes}
\begin{itemize}
\item fitting high-order global polynomials;
\item using observations far from the cutoff;
\item ignoring manipulation;
\item skipping covariate-balance checks;
\item treating a local effect as population-wide;
\item interpreting a discontinuity in treatment probability as a sharp design;
\item hiding bandwidth sensitivity.
\end{itemize}
\end{frame}

\begin{frame}{Synthesis: Identification at the Threshold}
\begin{itemize}
\item RDD exploits discontinuous treatment assignment at a known threshold.
\item Identification is local to units near the cutoff.
\item Continuity of untreated potential outcomes is the central design assumption.
\item Precise manipulation threatens local comparability.
\item Local linear estimation and principled bandwidth selection are preferred.
\item Fuzzy RDD is a local instrumental-variable design.
\item Density tests, covariate balance, placebo cutoffs, and bandwidth sensitivity are core diagnostics.
\end{itemize}
\end{frame}

\begin{frame}{Selected References}
\small
\begin{itemize}
\item Hahn, J., Todd, P., and van der Klaauw, W. (2001). Identification and Estimation of Treatment Effects with a Regression-Discontinuity Design. \textit{Econometrica}, 69(1), 201--209.
\item Imbens, G. W., and Lemieux, T. (2008). Regression Discontinuity Designs: A Guide to Practice. \textit{Journal of Econometrics}, 142(2), 615--635.
\item Lee, D. S., and Lemieux, T. (2010). Regression Discontinuity Designs in Economics. \textit{Journal of Economic Literature}, 48(2), 281--355.
\item Calonico, S., Cattaneo, M. D., and Titiunik, R. (2014). Robust Nonparametric Confidence Intervals for Regression-Discontinuity Designs. \textit{Econometrica}, 82(6), 2295--2326.
\item Cattaneo, M. D., Idrobo, N., and Titiunik, R. (2020). \textit{A Practical Introduction to Regression Discontinuity Designs}. Cambridge University Press.
\end{itemize}
\end{frame}

\contactslide

\end{document}
